\documentclass[12pt,a4paper]{article}

\usepackage[T2A]{fontenc}
\usepackage[utf8]{inputenc}
\usepackage[russian]{babel}
\usepackage{amsmath,amssymb,amsthm,mathtools}
\usepackage{enumitem}
\usepackage{geometry}
\usepackage{hyperref}
\usepackage{stmaryrd}
\usepackage{mathrsfs}
\usepackage{graphicx} % Для вставки картинок
\usepackage{float}    % Для жесткой фиксации картинок и схем на месте
\usepackage{tikz}     % Для рисования блок-схем
\usetikzlibrary{positioning, arrows.meta}

\geometry{margin=2.2cm}

\newtheorem{definition}{Определение}
\newtheorem{theorem}{Теорема}
\newtheorem{proposition}{Утверждение}
\newtheorem{lemma}{Лемма}
\newtheorem{corollary}{Следствие}
\theoremstyle{remark}
\newtheorem*{remark}{Замечание}
\newtheorem*{example}{Пример}

\title{Билет №22 \\ \small Архитектура современных компьютеров. Организации памяти и архитектура процессора современных вычислительных машин. Организация шин. Страничная и сегментная организация виртуальной памяти. Кэш-память. (СПбГУ)}
\author{Сериков Владислав Александрович}
\date{29 апреля 2026}

\begin{document}

\maketitle
\tableofcontents
\newpage

\section{Архитектура современных компьютеров}

\subsection{Понятие архитектуры компьютера}

\textbf{Архитектура компьютера} --- это совокупность принципов организации вычислительной машины, определяющих:
\begin{itemize}
    \item состав основных устройств;
    \item способы взаимодействия процессора, памяти и устройств ввода-вывода;
    \item систему команд процессора;
    \item организацию памяти;
    \item представление данных;
    \item механизмы адресации;
    \item способы повышения производительности: конвейеризация, кэширование, параллелизм и т.д.
\end{itemize}

В общем виде современный компьютер можно представить следующей схемой:

\[
\boxed{
\begin{array}{c}
\text{Процессор CPU} \\
\text{(ALU, CU, регистры, кэши)}
\end{array}}
\Longleftrightarrow
\boxed{
\begin{array}{c}
\text{Оперативная память RAM}
\end{array}}
\Longleftrightarrow
\boxed{
\begin{array}{c}
\text{Устройства ввода-вывода} \\
\text{диски, сеть, клавиатура, GPU}
\end{array}}
\]

Более детально:

\[
\begin{array}{cccccc}
 & & \boxed{\text{CPU}} & & & \\
 & & \begin{array}{c}
        \text{ALU} \\
        \text{CU} \\
        \text{Registers} \\
        \text{L1/L2 Cache}
      \end{array}
 & & & \\
 & & \updownarrow & & & \\
\boxed{\text{RAM}} & \Longleftrightarrow & \boxed{\text{System Bus / Interconnect}}
& \Longleftrightarrow & \boxed{\text{I/O Controllers}} & \Longleftrightarrow \boxed{\text{Devices}}
\end{array}
\]

Современные компьютеры обычно строятся на основе \textbf{модифицированной фон-неймановской архитектуры}, где:
\begin{itemize}
    \item программа и данные хранятся в общей памяти;
    \item процессор последовательно извлекает команды из памяти;
    \item команды декодируются и исполняются;
    \item результаты помещаются в регистры, память или устройства вывода.
\end{itemize}

Классическая фон-неймановская схема:

\[
\boxed{\text{Память}}
\Longleftrightarrow
\boxed{\text{Устройство управления}}
\Longleftrightarrow
\boxed{\text{АЛУ}}
\]

\[
\boxed{\text{Устройства ввода}}
\longrightarrow
\boxed{\text{Память / CPU}}
\longrightarrow
\boxed{\text{Устройства вывода}}
\]

Основная проблема такой архитектуры --- \textbf{фон-неймановское узкое место}: процессор обычно работает быстрее оперативной памяти, поэтому скорость выполнения программы часто ограничивается скоростью доставки команд и данных из памяти.

Для борьбы с этим применяются:
\begin{itemize}
    \item многоуровневая кэш-память;
    \item конвейеризация;
    \item предсказание переходов;
    \item внеочередное исполнение команд;
    \item многопроцессорность и многоядерность;
    \item SIMD-расширения;
    \item предварительная загрузка данных;
    \item виртуальная память.
\end{itemize}

\subsection{Классификация архитектур}

\subsubsection{Архитектура фон Неймана}

В архитектуре фон Неймана команды и данные хранятся в одной памяти и передаются по одной и той же шине.

\[
\boxed{\text{CPU}}
\Longleftrightarrow
\boxed{\text{Общая шина}}
\Longleftrightarrow
\boxed{\text{Память команд и данных}}
\]

Преимущества:
\begin{itemize}
    \item простота организации;
    \item универсальность;
    \item возможность хранить программы как данные.
\end{itemize}

Недостатки:
\begin{itemize}
    \item команды и данные конкурируют за одну шину;
    \item возникает ограничение пропускной способности памяти;
    \item сложнее обеспечить параллельную выборку команд и данных.
\end{itemize}

\subsubsection{Гарвардская архитектура}

В гарвардской архитектуре память команд и память данных разделены.

\[
\boxed{\text{CPU}}
\begin{array}{c}
\Longleftrightarrow \boxed{\text{Память команд}} \\
\Longleftrightarrow \boxed{\text{Память данных}}
\end{array}
\]

Преимущества:
\begin{itemize}
    \item можно одновременно выбирать команду и читать/записывать данные;
    \item выше производительность;
    \item удобно для микроконтроллеров и цифровых сигнальных процессоров.
\end{itemize}

Недостатки:
\begin{itemize}
    \item менее гибкая организация памяти;
    \item сложнее программная модель.
\end{itemize}

В современных процессорах часто используется \textbf{модифицированная гарвардская архитектура}:
\begin{itemize}
    \item на уровне оперативной памяти команды и данные находятся в общем адресном пространстве;
    \item на уровне кэша первого уровня часто разделяются кэш команд и кэш данных.
\end{itemize}

Например:

\[
\begin{array}{ccc}
 & \boxed{\text{CPU Core}} & \\
 & \begin{array}{cc}
        \boxed{\text{L1 I-cache}} & \boxed{\text{L1 D-cache}}
   \end{array} & \\
 & \downarrow & \\
 & \boxed{\text{L2 Cache}} & \\
 & \downarrow & \\
 & \boxed{\text{L3 Cache}} & \\
 & \downarrow & \\
 & \boxed{\text{RAM}} &
\end{array}
\]

\subsection{Основные компоненты современного компьютера}

Современная вычислительная машина включает:

\begin{enumerate}
    \item \textbf{Центральный процессор} --- выполняет команды программы.
    \item \textbf{Оперативную память} --- хранит исполняемые программы и данные.
    \item \textbf{Постоянную память} --- SSD, HDD, Flash.
    \item \textbf{Системные шины и межсоединения} --- обеспечивают обмен данными.
    \item \textbf{Контроллеры ввода-вывода} --- управляют периферийными устройствами.
    \item \textbf{Графический процессор} --- выполняет графические и параллельные вычисления.
    \item \textbf{Сетевые адаптеры} --- обеспечивают сетевое взаимодействие.
\end{enumerate}

Упрощённая структура:

\[
\begin{array}{ccccc}
\boxed{\text{CPU}}
& \Longleftrightarrow &
\boxed{\text{Memory Controller}}
& \Longleftrightarrow &
\boxed{\text{RAM}}
\\[1em]
\updownarrow & & & & \\
\boxed{\text{PCI Express}}
& \Longleftrightarrow &
\boxed{\text{GPU, SSD, NIC}}
& &
\end{array}
\]

В старых компьютерах часто использовалась схема с северным и южным мостами:

\[
\begin{array}{ccccc}
 & \boxed{\text{CPU}} & & & \\
 & \updownarrow & & & \\
 & \boxed{\text{Northbridge}} & \Longleftrightarrow & \boxed{\text{RAM}} & \\
 & \updownarrow & & & \\
 & \boxed{\text{Southbridge}} & \Longleftrightarrow & \boxed{\text{I/O devices}} &
\end{array}
\]

В современных системах контроллер памяти обычно встроен непосредственно в процессор:

\[
\boxed{\text{CPU + Memory Controller}}
\Longleftrightarrow
\boxed{\text{RAM}}
\]

Это уменьшает задержки доступа к памяти.

\section{Архитектура процессора современных вычислительных машин}

\subsection{Основные блоки процессора}

\textbf{Центральный процессор} --- устройство, выполняющее машинные команды.

Основные компоненты процессора:

\begin{itemize}
    \item \textbf{Устройство управления} --- извлекает, декодирует и организует выполнение команд.
    \item \textbf{Арифметико-логическое устройство} --- выполняет арифметические и логические операции.
    \item \textbf{Регистры} --- сверхбыстрая внутренняя память процессора.
    \item \textbf{Блок работы с памятью} --- выполняет загрузку и сохранение данных.
    \item \textbf{Блок предсказания переходов}.
    \item \textbf{Конвейер команд}.
    \item \textbf{Кэш-память}.
    \item \textbf{Блок управления виртуальной памятью}, или MMU.
    \item \textbf{Векторные/SIMD-блоки}.
\end{itemize}

Схема ядра процессора:

\[
\begin{array}{c}
\boxed{
\begin{array}{c}
\text{CPU Core} \\[0.5em]
\boxed{\text{Fetch Unit}} \\
\boxed{\text{Decode Unit}} \\
\boxed{\text{Control Unit}} \\
\boxed{\text{ALU / FPU / SIMD}} \\
\boxed{\text{Registers}} \\
\boxed{\text{Load/Store Unit}} \\
\boxed{\text{MMU + TLB}} \\
\boxed{\text{L1 Cache}}
\end{array}}
\end{array}
\]

\subsection{Регистры процессора}

\textbf{Регистры} --- это небольшие ячейки памяти внутри процессора, доступ к которым выполняется быстрее, чем к кэшу и оперативной памяти.

Типы регистров:
\begin{itemize}
    \item \textbf{регистры общего назначения};
    \item \textbf{счётчик команд} PC/IP;
    \item \textbf{регистр состояния} FLAGS/PSW;
    \item \textbf{указатель стека} SP;
    \item \textbf{регистры адреса};
    \item \textbf{векторные регистры};
    \item \textbf{управляющие регистры}.
\end{itemize}

Например, команда:

\[
C = A + B
\]

на машинном уровне может выполняться так:

\[
\begin{array}{l}
\texttt{LOAD R1, [A]} \\
\texttt{LOAD R2, [B]} \\
\texttt{ADD R3, R1, R2} \\
\texttt{STORE [C], R3}
\end{array}
\]

Здесь:
\begin{itemize}
    \item \texttt{R1}, \texttt{R2}, \texttt{R3} --- регистры;
    \item \texttt{LOAD} загружает данные из памяти;
    \item \texttt{ADD} выполняет сложение;
    \item \texttt{STORE} сохраняет результат в память.
\end{itemize}

\subsection{Арифметико-логическое устройство}

\textbf{АЛУ} выполняет операции:
\begin{itemize}
    \item сложение;
    \item вычитание;
    \item умножение;
    \item деление;
    \item логическое И;
    \item логическое ИЛИ;
    \item исключающее ИЛИ;
    \item сдвиги;
    \item сравнения.
\end{itemize}

Схематично:

\[
\boxed{\text{Операнд 1}}
\longrightarrow
\boxed{\text{ALU}}
\longleftarrow
\boxed{\text{Операнд 2}}
\]

\[
\boxed{\text{ALU}}
\longrightarrow
\boxed{\text{Результат}}
\]

\subsection{Устройство управления}

\textbf{Устройство управления} организует цикл выполнения команды:

\[
\boxed{\text{Fetch}}
\rightarrow
\boxed{\text{Decode}}
\rightarrow
\boxed{\text{Execute}}
\rightarrow
\boxed{\text{Memory}}
\rightarrow
\boxed{\text{Write Back}}
\]

Этапы:
\begin{enumerate}
    \item \textbf{Fetch} --- выборка команды из памяти.
    \item \textbf{Decode} --- декодирование команды.
    \item \textbf{Execute} --- выполнение операции.
    \item \textbf{Memory Access} --- обращение к памяти, если нужно.
    \item \textbf{Write Back} --- запись результата.
\end{enumerate}

Пример:

\[
\texttt{ADD R1, R2, R3}
\]

означает:

\[
R1 := R2 + R3.
\]

Эта команда:
\begin{itemize}
    \item выбирается из памяти;
    \item декодируется;
    \item операнды \texttt{R2} и \texttt{R3} читаются из регистров;
    \item АЛУ складывает их;
    \item результат записывается в \texttt{R1}.
\end{itemize}

\subsection{Конвейеризация}

\textbf{Конвейеризация} --- способ повышения производительности процессора, при котором разные стадии выполнения нескольких команд выполняются параллельно.

Без конвейера:

\[
\begin{array}{c|ccccc}
\text{Такт} & 1 & 2 & 3 & 4 & 5 \\ \hline
I_1 & F & D & E & M & W \\
I_2 &  &  &  &  & F
\end{array}
\]

С конвейером:

\[
\begin{array}{c|ccccccccc}
\text{Такт} & 1 & 2 & 3 & 4 & 5 & 6 & 7 & 8 & 9 \\ \hline
I_1 & F & D & E & M & W & & & & \\
I_2 &   & F & D & E & M & W & & & \\
I_3 &   &   & F & D & E & M & W & & \\
I_4 &   &   &   & F & D & E & M & W & \\
I_5 &   &   &   &   & F & D & E & M & W
\end{array}
\]

Где:
\[
F = Fetch,\quad D = Decode,\quad E = Execute,\quad M = Memory,\quad W = WriteBack.
\]

Если одна команда выполняется за 5 стадий, то без конвейера 5 команд заняли бы:

\[
5 \cdot 5 = 25 \text{ тактов}.
\]

С конвейером:

\[
5 + (5 - 1) = 9 \text{ тактов}.
\]

Однако существуют конфликты:
\begin{itemize}
    \item \textbf{структурные} --- нескольким командам нужен один и тот же ресурс;
    \item \textbf{по данным} --- команда зависит от результата предыдущей;
    \item \textbf{управляющие} --- переходы и ветвления.
\end{itemize}

Пример конфликта по данным:

\[
\begin{array}{l}
\texttt{ADD R1, R2, R3} \\
\texttt{SUB R4, R1, R5}
\end{array}
\]

Вторая команда использует \texttt{R1}, который ещё не успел быть записан первой командой.

\subsection{Суперскалярность}

\textbf{Суперскалярный процессор} может запускать на выполнение несколько команд за один такт.

Например, если процессор имеет два исполнительных устройства:

\[
\begin{array}{c|cc}
\text{Такт} & \text{ALU 1} & \text{ALU 2} \\ \hline
1 & I_1 & I_2 \\
2 & I_3 & I_4 \\
3 & I_5 & I_6
\end{array}
\]

Это повышает производительность, если команды независимы.

\subsection{Внеочередное исполнение}

\textbf{Внеочередное исполнение} означает, что процессор может выполнять команды не строго в порядке их расположения в программе, если это не нарушает зависимостей по данным.

Пример программы:

\[
\begin{array}{l}
I_1: \texttt{LOAD R1, [A]} \\
I_2: \texttt{ADD R2, R1, R3} \\
I_3: \texttt{MUL R4, R5, R6}
\end{array}
\]

Если загрузка \texttt{R1} задерживается, процессор может временно выполнить независимую команду:

\[
I_3: \texttt{MUL R4, R5, R6}.
\]

\subsection{Предсказание переходов}

Команды условного перехода нарушают работу конвейера, потому что процессору заранее неизвестно, какая команда будет следующей.

Пример:

\[
\begin{array}{l}
\texttt{CMP R1, 0} \\
\texttt{JZ label}
\end{array}
\]

Если процессор ждёт результата сравнения, конвейер простаивает. Поэтому используется \textbf{предсказание переходов}: процессор предполагает, будет ли переход выполнен, и заранее загружает соответствующие команды.

Если предсказание верное, производительность растёт. Если неверное --- конвейер очищается, и возникает штраф по тактам.

\subsection{RISC и CISC}

\textbf{CISC} --- Complex Instruction Set Computer, архитектура с большим и сложным набором команд.

Примеры: x86, x86-64.

Особенности:
\begin{itemize}
    \item сложные инструкции;
    \item команды переменной длины;
    \item много режимов адресации;
    \item одна команда может выполнять сложное действие.
\end{itemize}

Например:

\[
\texttt{ADD [A], R1}
\]

может означать: взять значение из памяти по адресу \texttt{A}, прибавить \texttt{R1}, результат записать обратно в память.

\textbf{RISC} --- Reduced Instruction Set Computer, архитектура с упрощённым набором команд.

Примеры: ARM, RISC-V, MIPS.

Особенности:
\begin{itemize}
    \item простые команды;
    \item фиксированная длина команд;
    \item операции с памятью обычно только \texttt{LOAD}/\texttt{STORE};
    \item удобно для конвейеризации.
\end{itemize}

Пример RISC-подхода:

\[
\begin{array}{l}
\texttt{LOAD R2, [A]} \\
\texttt{ADD R2, R2, R1} \\
\texttt{STORE [A], R2}
\end{array}
\]

Современные процессоры x86 внешне являются CISC, но внутри часто преобразуют сложные команды в более простые микрооперации, близкие по духу к RISC.

\section{Организация памяти современных вычислительных машин}

\subsection{Иерархия памяти}

Память компьютера организована иерархически.

\[
\begin{array}{c}
\boxed{\text{Регистры}} \\
\downarrow \\
\boxed{\text{Кэш L1}} \\
\downarrow \\
\boxed{\text{Кэш L2}} \\
\downarrow \\
\boxed{\text{Кэш L3}} \\
\downarrow \\
\boxed{\text{Оперативная память RAM}} \\
\downarrow \\
\boxed{\text{SSD / HDD}} \\
\downarrow \\
\boxed{\text{Архивная память}}
\end{array}
\]

Чем выше уровень памяти:
\begin{itemize}
    \item тем меньше объём;
    \item тем выше скорость;
    \item тем дороже стоимость одного байта.
\end{itemize}

Чем ниже уровень:
\begin{itemize}
    \item тем больше объём;
    \item тем ниже скорость;
    \item тем дешевле хранение.
\end{itemize}

Таблица иерархии памяти:

\[
\begin{array}{|c|c|c|c|}
\hline
\text{Уровень} & \text{Примерный объём} & \text{Скорость} & \text{Назначение} \\
\hline
\text{Регистры} & \text{байты--КБ} & \text{максимальная} & \text{операнды команд} \\
\hline
\text{L1 Cache} & \text{десятки КБ} & \text{очень высокая} & \text{ближайший кэш} \\
\hline
\text{L2 Cache} & \text{сотни КБ--МБ} & \text{высокая} & \text{кэш ядра} \\
\hline
\text{L3 Cache} & \text{МБ--десятки МБ} & \text{средняя} & \text{общий кэш} \\
\hline
\text{RAM} & \text{ГБ--ТБ} & \text{ниже кэша} & \text{программы и данные} \\
\hline
\text{SSD/HDD} & \text{сотни ГБ--ТБ} & \text{медленная} & \text{долговременное хранение} \\
\hline
\end{array}
\]

\subsection{Принцип локальности}

Эффективность иерархии памяти основана на \textbf{принципе локальности}.

Различают:

\begin{enumerate}
    \item \textbf{Временную локальность}: если данные использовались недавно, вероятно, они скоро понадобятся снова.
    \item \textbf{Пространственную локальность}: если использовался некоторый адрес, вероятно, скоро будут использоваться соседние адреса.
\end{enumerate}

Пример временной локальности:

\[
\begin{array}{l}
\texttt{for (i = 0; i < 1000000; i++)} \\
\quad \texttt{sum += x;}
\end{array}
\]

Переменная \texttt{x} многократно используется.

Пример пространственной локальности:

\[
\begin{array}{l}
\texttt{for (i = 0; i < n; i++)} \\
\quad \texttt{sum += a[i];}
\end{array}
\]

Элементы массива \texttt{a[i]} расположены в памяти подряд.

\subsection{Физическая и виртуальная память}

\textbf{Физическая память} --- реальная оперативная память компьютера, представленная микросхемами RAM.

\textbf{Виртуальная память} --- абстракция, предоставляемая операционной системой каждому процессу. Каждый процесс как будто имеет своё собственное непрерывное адресное пространство.

Например:

\[
\begin{array}{c}
\boxed{\text{Процесс 1: виртуальная память}} \\
0x00000000 \ldots 0xFFFFFFFF
\end{array}
\]

\[
\begin{array}{c}
\boxed{\text{Процесс 2: виртуальная память}} \\
0x00000000 \ldots 0xFFFFFFFF
\end{array}
\]

Хотя виртуальные адреса могут совпадать, они отображаются на разные физические адреса.

\[
\begin{array}{ccc}
\text{Процесс 1: VA } 0x1000 & \longrightarrow & \text{PA } 0xA000 \\
\text{Процесс 2: VA } 0x1000 & \longrightarrow & \text{PA } 0xB000
\end{array}
\]

Где:
\begin{itemize}
    \item VA --- virtual address, виртуальный адрес;
    \item PA --- physical address, физический адрес.
\end{itemize}

Преимущества виртуальной памяти:
\begin{itemize}
    \item изоляция процессов;
    \item защита памяти;
    \item удобная модель программирования;
    \item возможность использовать больше памяти, чем есть физически;
    \item поддержка отображения файлов в память;
    \item разделяемая память между процессами.
\end{itemize}

\section{Организация шин}

\subsection{Понятие шины}

\textbf{Шина} --- это система линий связи, по которым передаются данные, адреса и управляющие сигналы между устройствами компьютера.

Классически выделяют три типа шин:

\begin{enumerate}
    \item \textbf{Шина данных} --- передаёт сами данные.
    \item \textbf{Шина адреса} --- указывает адрес ячейки памяти или устройства.
    \item \textbf{Шина управления} --- передаёт сигналы чтения, записи, прерывания, синхронизации.
\end{enumerate}

Схема:

\[
\boxed{\text{CPU}}
\begin{array}{c}
\Longleftrightarrow \boxed{\text{Шина данных}} \\
\Longrightarrow \boxed{\text{Шина адреса}} \\
\Longleftrightarrow \boxed{\text{Шина управления}}
\end{array}
\boxed{\text{Memory / I/O}}
\]

\subsection{Шина данных}

\textbf{Шина данных} определяет, сколько бит может быть передано за один цикл.

Например:
\begin{itemize}
    \item 32-битная шина передаёт 32 бита за один такт;
    \item 64-битная шина передаёт 64 бита за один такт.
\end{itemize}

Если ширина шины данных равна $w$ бит, а частота передачи $f$, то теоретическая пропускная способность:

\[
B = w \cdot f.
\]

Если шина 64-битная и работает на частоте $1$ ГГц:

\[
B = 64 \cdot 10^9 \text{ бит/с} = 8 \cdot 10^9 \text{ байт/с} = 8 \text{ ГБ/с}.
\]

\subsection{Шина адреса}

\textbf{Шина адреса} задаёт, сколько различных адресов может быть сформировано.

Если ширина адресной шины равна $n$ бит, то можно адресовать:

\[
2^n
\]

различных адресов.

Если адресуется каждый байт, то максимальный объём адресуемой памяти:

\[
M = 2^n \text{ байт}.
\]

Например, для 32-битной адресной шины:

\[
M = 2^{32} \text{ байт} = 4 \text{ ГБ}.
\]

Для 64-битной:

\[
M = 2^{64} \text{ байт}.
\]

На практике современные процессоры могут использовать не все 64 бита физического адреса.

\subsection{Шина управления}

\textbf{Шина управления} содержит сигналы:
\begin{itemize}
    \item чтение из памяти;
    \item запись в память;
    \item подтверждение готовности;
    \item прерывания;
    \item запросы прямого доступа к памяти;
    \item тактовые сигналы;
    \item сигналы сброса.
\end{itemize}

Пример операции чтения из памяти:

\[
\begin{array}{rcl}
1. & \text{CPU} & \text{выставляет адрес на адресную шину;} \\
2. & \text{CPU} & \text{подаёт сигнал READ;} \\
3. & \text{Memory} & \text{помещает данные на шину данных;} \\
4. & \text{CPU} & \text{считывает данные.}
\end{array}
\]

\subsection{Системная шина и современные межсоединения}

В классической архитектуре используется \textbf{системная шина}, соединяющая CPU, память и устройства ввода-вывода.

\[
\begin{array}{ccc}
\boxed{\text{CPU}} & & \boxed{\text{RAM}} \\
 & \diagup \quad \diagdown & \\
 & \boxed{\text{System Bus}} & \\
 & \diagdown \quad \diagup & \\
\boxed{\text{I/O}} & & \boxed{\text{DMA}}
\end{array}
\]

Однако единая общая шина плохо масштабируется: если много устройств одновременно хотят передавать данные, возникает конкуренция.

Поэтому современные системы используют:
\begin{itemize}
    \item \textbf{точка-точка соединения};
    \item \textbf{PCI Express};
    \item \textbf{встроенный контроллер памяти};
    \item \textbf{многоуровневые межсоединения};
    \item \textbf{NoC} --- Network-on-Chip в многоядерных процессорах.
\end{itemize}

\subsection{PCI Express}

\textbf{PCI Express} --- современная высокоскоростная последовательная шина для подключения видеокарт, SSD, сетевых карт и других устройств.

PCIe использует линии, или lanes. Одна линия состоит из двух дифференциальных пар:
\begin{itemize}
    \item одна пара для передачи;
    \item одна пара для приёма.
\end{itemize}

Конфигурации:

\[
x1,\quad x4,\quad x8,\quad x16.
\]

Например, видеокарта обычно подключается через PCIe $x16$.

Схема:

\[
\boxed{\text{CPU / Chipset}}
\Longleftrightarrow
\boxed{\text{PCIe x16}}
\Longleftrightarrow
\boxed{\text{GPU}}
\]

\subsection{DMA}

\textbf{DMA} --- Direct Memory Access, прямой доступ к памяти.

DMA позволяет устройствам ввода-вывода обмениваться данными с оперативной памятью без постоянного участия процессора.

Без DMA:

\[
\boxed{\text{Disk}} \rightarrow \boxed{\text{CPU}} \rightarrow \boxed{\text{RAM}}
\]

С DMA:

\[
\boxed{\text{Disk}} \rightarrow \boxed{\text{RAM}}
\]

Процессор только настраивает передачу:
\begin{enumerate}
    \item указывает адрес буфера в памяти;
    \item задаёт размер блока;
    \item запускает контроллер DMA;
    \item после завершения получает прерывание.
\end{enumerate}

Преимущество:
\[
\text{CPU освобождается от рутинной пересылки данных.}
\]

\section{Страничная организация виртуальной памяти}

\subsection{Основная идея}

\textbf{Страничная организация памяти} --- способ управления виртуальной памятью, при котором:
\begin{itemize}
    \item виртуальное адресное пространство делится на блоки фиксированного размера --- страницы;
    \item физическая память делится на блоки такого же размера --- страничные кадры;
    \item таблица страниц задаёт соответствие между виртуальными страницами и физическими кадрами.
\end{itemize}

\[
\boxed{\text{Виртуальная память}}
=
\boxed{\text{страница 0}}
\boxed{\text{страница 1}}
\boxed{\text{страница 2}}
\cdots
\]

\[
\boxed{\text{Физическая память}}
=
\boxed{\text{кадр 0}}
\boxed{\text{кадр 1}}
\boxed{\text{кадр 2}}
\cdots
\]

Отображение:

\[
\text{Virtual Page Number} \longrightarrow \text{Physical Frame Number}.
\]

\subsection{Структура виртуального адреса}

Виртуальный адрес делится на две части:

\[
\boxed{\text{Номер виртуальной страницы}}
\quad
\boxed{\text{Смещение внутри страницы}}
\]

То есть:

\[
VA = (VPN, offset).
\]

Если размер страницы равен $2^k$ байт, то младшие $k$ бит адреса являются смещением.

Например, если размер страницы $4$ КБ:

\[
4 \text{ КБ} = 4096 = 2^{12}.
\]

Значит:

\[
offset = 12 \text{ бит}.
\]

Для 32-битного виртуального адреса:

\[
\underbrace{31 \ldots 12}_{20 \text{ бит: VPN}}
\quad
\underbrace{11 \ldots 0}_{12 \text{ бит: offset}}
\]

Схема:

\[
\boxed{\text{VPN: 20 бит}}
\boxed{\text{offset: 12 бит}}
\]

Количество виртуальных страниц:

\[
2^{20}.
\]

\subsection{Преобразование адреса}

Пусть:
\begin{itemize}
    \item виртуальный адрес состоит из номера страницы $p$ и смещения $d$;
    \item таблица страниц содержит номер физического кадра $f$.
\end{itemize}

Тогда физический адрес:

\[
PA = f \cdot page\_size + d.
\]

Или в двоичном виде:

\[
PA = \boxed{\text{номер физического кадра}}
\boxed{\text{смещение}}.
\]

Схема преобразования:

\[
\boxed{\text{Virtual Address}}
\rightarrow
\boxed{\text{Page Number}}
\rightarrow
\boxed{\text{Page Table}}
\rightarrow
\boxed{\text{Frame Number}}
\rightarrow
\boxed{\text{Physical Address}}
\]

Более подробно:

\[
\begin{array}{ccc}
\boxed{VPN} & \boxed{offset} & \\
\downarrow & & \\
\boxed{\text{Page Table}} & & \\
\downarrow & & \\
\boxed{PFN} & \boxed{offset} & \longrightarrow \boxed{PA}
\end{array}
\]

\subsection{Пример преобразования адреса}

Пусть:
\begin{itemize}
    \item размер страницы $= 4$ КБ $= 4096$ байт;
    \item виртуальный адрес $VA = 0x12345$;
    \item таблица страниц содержит:
    \[
    VPN = 0x12 \longrightarrow PFN = 0xAB.
    \]
\end{itemize}

Так как размер страницы $2^{12}$, смещение занимает 12 бит.

\[
VA = 0x12345.
\]

Разделим:

\[
VPN = 0x12,
\quad
offset = 0x345.
\]

По таблице страниц:

\[
VPN = 0x12 \longrightarrow PFN = 0xAB.
\]

Тогда физический адрес:

\[
PA = 0xAB345.
\]

\subsection{Запись таблицы страниц}

Каждая запись таблицы страниц называется \textbf{PTE} --- Page Table Entry.

Она обычно содержит:
\begin{itemize}
    \item номер физического кадра;
    \item бит присутствия страницы в памяти;
    \item бит разрешения чтения;
    \item бит разрешения записи;
    \item бит разрешения исполнения;
    \item бит модификации;
    \item бит обращения;
    \item признаки режима пользователя/ядра;
    \item признаки кэширования.
\end{itemize}

Пример:

\[
\begin{array}{|c|c|c|c|c|}
\hline
VPN & PFN & Present & R/W & User/Supervisor \\
\hline
0x12 & 0xAB & 1 & R/W & User \\
\hline
0x13 & - & 0 & - & - \\
\hline
\end{array}
\]

Если \texttt{Present = 0}, страница отсутствует в оперативной памяти. При обращении возникает \textbf{page fault}.

\subsection{Page fault}

\textbf{Page fault} --- исключительная ситуация, возникающая при обращении к странице, которой нет в физической памяти, либо доступ к которой запрещён.

Возможные причины:
\begin{itemize}
    \item страница выгружена на диск;
    \item процесс обращается к невыделенной памяти;
    \item нарушение прав доступа;
    \item попытка записи в страницу только для чтения;
    \item попытка исполнения неисполняемой страницы.
\end{itemize}

Обработка page fault:

\[
\begin{array}{rcl}
1. & \text{CPU} & \text{генерирует исключение page fault;} \\
2. & \text{ОС} & \text{проверяет корректность адреса;} \\
3. & \text{ОС} & \text{если страница допустима, загружает её в RAM;} \\
4. & \text{ОС} & \text{обновляет таблицу страниц;} \\
5. & \text{CPU} & \text{повторяет команду.}
\end{array}
\]

Если доступ недопустим, процесс обычно завершается с ошибкой, например \texttt{segmentation fault}.

\subsection{Многоуровневые таблицы страниц}

Одноуровневая таблица страниц может быть слишком большой.

Например, для 32-битного адресного пространства и страниц по 4 КБ:

\[
2^{32} / 2^{12} = 2^{20}
\]

страниц.

Если каждая запись таблицы страниц занимает 4 байта:

\[
2^{20} \cdot 4 = 4 \text{ МБ}
\]

на один процесс.

Для 64-битных адресов одноуровневая таблица была бы слишком огромной.

Поэтому применяются многоуровневые таблицы страниц.

Например, двухуровневая схема для 32-битного адреса:

\[
\boxed{\text{Page Directory Index}}
\boxed{\text{Page Table Index}}
\boxed{\text{Offset}}
\]

Если страница 4 КБ, offset = 12 бит. Оставшиеся 20 бит можно разбить:

\[
10 + 10 + 12.
\]

\[
\underbrace{31 \ldots 22}_{10 \text{ бит: каталог}}
\quad
\underbrace{21 \ldots 12}_{10 \text{ бит: таблица}}
\quad
\underbrace{11 \ldots 0}_{12 \text{ бит: смещение}}
\]

Схема:

\[
\boxed{\text{VA}}
\rightarrow
\boxed{\text{Page Directory}}
\rightarrow
\boxed{\text{Page Table}}
\rightarrow
\boxed{\text{Physical Frame}}
\rightarrow
\boxed{\text{PA}}
\]

\subsection{TLB}

\textbf{TLB} --- Translation Lookaside Buffer, ассоциативный кэш трансляций виртуальных адресов в физические.

Без TLB для каждого обращения к памяти нужно сначала обратиться к таблице страниц, а затем к самой памяти. Это сильно замедлило бы работу.

С TLB:

\[
\boxed{VPN}
\rightarrow
\boxed{\text{TLB}}
\rightarrow
\boxed{PFN}
\]

Если запись найдена, происходит \textbf{TLB hit}. Если нет --- \textbf{TLB miss}, и процессор обращается к таблице страниц.

Схема:

\[
\begin{array}{ccc}
\boxed{\text{Virtual Address}} & \longrightarrow & \boxed{\text{TLB}} \\
 & \searrow & \downarrow \\
 & & \boxed{\text{Page Table}} \\
 & & \downarrow \\
 & & \boxed{\text{Physical Address}}
\end{array}
\]

\subsection{Достоинства и недостатки страничной организации}

Достоинства:
\begin{itemize}
    \item нет внешней фрагментации;
    \item удобно выделять память блоками фиксированного размера;
    \item легко реализовать виртуальную память;
    \item удобно защищать память;
    \item можно разделять страницы между процессами;
    \item можно отображать файлы в память.
\end{itemize}

Недостатки:
\begin{itemize}
    \item есть внутренняя фрагментация;
    \item нужны таблицы страниц;
    \item требуется аппаратная поддержка MMU;
    \item возможны накладные расходы на page fault;
    \item нужен TLB для высокой производительности.
\end{itemize}

\section{Сегментная организация виртуальной памяти}

\subsection{Основная идея}

\textbf{Сегментная организация памяти} делит виртуальное адресное пространство не на равные страницы, а на логические области переменного размера --- сегменты.

Типичные сегменты:
\begin{itemize}
    \item сегмент кода;
    \item сегмент данных;
    \item сегмент стека;
    \item сегмент кучи;
    \item сегмент библиотек.
\end{itemize}

Схема виртуального адресного пространства процесса:

\[
\begin{array}{c}
\boxed{\text{Код}} \\
\boxed{\text{Глобальные данные}} \\
\boxed{\text{Куча}} \\
\boxed{\text{Свободная область}} \\
\boxed{\text{Стек}}
\end{array}
\]

При сегментации виртуальный адрес имеет вид:

\[
VA = (s, d),
\]

где:
\begin{itemize}
    \item $s$ --- номер сегмента;
    \item $d$ --- смещение внутри сегмента.
\end{itemize}

\[
\boxed{\text{Номер сегмента}}
\quad
\boxed{\text{Смещение}}
\]

\subsection{Таблица сегментов}

Каждый процесс имеет таблицу сегментов.

Запись таблицы сегментов содержит:
\begin{itemize}
    \item базовый физический адрес сегмента;
    \item длину сегмента;
    \item права доступа;
    \item признаки присутствия;
    \item признаки защиты.
\end{itemize}

Пример:

\[
\begin{array}{|c|c|c|c|}
\hline
\text{Сегмент} & \text{База} & \text{Длина} & \text{Права} \\
\hline
0: \text{Код} & 0x10000 & 0x4000 & R-X \\
\hline
1: \text{Данные} & 0x20000 & 0x3000 & RW- \\
\hline
2: \text{Стек} & 0x30000 & 0x2000 & RW- \\
\hline
\end{array}
\]

\subsection{Преобразование адреса при сегментации}

Пусть виртуальный адрес:

\[
VA = (s, d).
\]

Процессор:
\begin{enumerate}
    \item по номеру $s$ находит запись в таблице сегментов;
    \item проверяет, что $d < limit_s$;
    \item проверяет права доступа;
    \item вычисляет физический адрес:
    \[
    PA = base_s + d.
    \]
\end{enumerate}

Схема:

\[
\boxed{s}
\rightarrow
\boxed{\text{Segment Table}}
\rightarrow
\boxed{base, limit, rights}
\]

\[
\boxed{d}
\rightarrow
\boxed{d < limit?}
\rightarrow
\boxed{PA = base + d}
\]

Пример:

Пусть:
\[
s = 1,\quad d = 0x120.
\]

Из таблицы:

\[
base_1 = 0x20000,\quad limit_1 = 0x3000.
\]

Проверяем:

\[
0x120 < 0x3000.
\]

Тогда:

\[
PA = 0x20000 + 0x120 = 0x20120.
\]

\subsection{Достоинства сегментации}

\begin{itemize}
    \item соответствует логической структуре программы;
    \item удобно задавать разные права для разных частей программы;
    \item легко разделять сегменты между процессами;
    \item сегмент может расти независимо от других;
    \item удобно защищать код от записи, стек от исполнения и т.д.
\end{itemize}

Например:
\[
\text{код: только чтение и исполнение},
\]
\[
\text{данные: чтение и запись},
\]
\[
\text{стек: чтение и запись, но не исполнение}.
\]

\subsection{Недостатки сегментации}

\begin{itemize}
    \item сегменты имеют переменный размер;
    \item возникает внешняя фрагментация;
    \item сложнее выделять и освобождать память;
    \item требуется уплотнение памяти;
    \item сложнее аппаратная реализация.
\end{itemize}

\textbf{Внешняя фрагментация} --- ситуация, когда суммарно свободной памяти достаточно, но она разбита на мелкие куски, и крупный сегмент разместить невозможно.

Пример:

\[
\boxed{\text{занято 10}}
\boxed{\text{свободно 5}}
\boxed{\text{занято 20}}
\boxed{\text{свободно 7}}
\boxed{\text{занято 15}}
\]

Свободно всего:

\[
5 + 7 = 12,
\]

но сегмент размера 10 можно разместить не всегда, если нет подходящего непрерывного участка.

\section{Сравнение страничной и сегментной организации}

\[
\begin{array}{|c|c|c|}
\hline
\text{Признак} & \text{Страничная организация} & \text{Сегментная организация} \\
\hline
\text{Размер блоков} & \text{фиксированный} & \text{переменный} \\
\hline
\text{Единица деления} & \text{страница} & \text{логический сегмент} \\
\hline
\text{Фрагментация} & \text{внутренняя} & \text{внешняя} \\
\hline
\text{Адрес} & (VPN, offset) & (segment, offset) \\
\hline
\text{Таблица} & \text{таблица страниц} & \text{таблица сегментов} \\
\hline
\text{Соответствие программе} & \text{слабое} & \text{сильное} \\
\hline
\text{Защита} & \text{на уровне страниц} & \text{на уровне сегментов} \\
\hline
\text{Удобство выделения памяти} & \text{выше} & \text{ниже} \\
\hline
\end{array}
\]

В современных операционных системах чаще используется страничная виртуальная память, иногда с элементами сегментации. Например, в x86 исторически поддерживается сегментация, но в 64-битном режиме основная роль принадлежит страничной адресации.

\section{Сегментно-страничная организация памяти}

\subsection{Основная идея}

В сегментно-страничной организации:
\begin{itemize}
    \item адресное пространство делится на логические сегменты;
    \item каждый сегмент, в свою очередь, делится на страницы;
    \item физическая память делится на страничные кадры.
\end{itemize}

Виртуальный адрес:

\[
VA = (s, p, d),
\]

где:
\begin{itemize}
    \item $s$ --- номер сегмента;
    \item $p$ --- номер страницы внутри сегмента;
    \item $d$ --- смещение внутри страницы.
\end{itemize}

\[
\boxed{\text{Segment}}
\boxed{\text{Page}}
\boxed{\text{Offset}}
\]

Схема:

\[
\boxed{s}
\rightarrow
\boxed{\text{Segment Table}}
\rightarrow
\boxed{\text{Page Table of Segment}}
\]

\[
\boxed{p}
\rightarrow
\boxed{\text{Page Table}}
\rightarrow
\boxed{PFN}
\]

\[
\boxed{PFN}
\boxed{d}
\rightarrow
\boxed{PA}
\]

Преимущества:
\begin{itemize}
    \item логическая структура сегментов;
    \item отсутствие внешней фрагментации благодаря страницам;
    \item гибкая защита;
    \item возможность выгружать части сегмента.
\end{itemize}

Недостатки:
\begin{itemize}
    \item сложнее аппаратная реализация;
    \item больше таблиц;
    \item больше накладные расходы на трансляцию адресов.
\end{itemize}

\section{Кэш-память}

\subsection{Назначение кэш-памяти}

\textbf{Кэш-память} --- это быстрая память небольшого объёма, расположенная между процессором и основной памятью. Она хранит копии часто используемых данных и команд.

Главная цель:

\[
\text{уменьшить среднее время доступа к памяти}.
\]

Проблема:

\[
\text{CPU работает намного быстрее RAM}.
\]

Поэтому без кэша процессор часто простаивал бы в ожидании данных.

Схема:

\[
\boxed{\text{CPU}}
\Longleftrightarrow
\boxed{\text{Cache}}
\Longleftrightarrow
\boxed{\text{RAM}}
\]

Многоуровневая схема:

\[
\boxed{\text{CPU}}
\Longleftrightarrow
\boxed{\text{L1}}
\Longleftrightarrow
\boxed{\text{L2}}
\Longleftrightarrow
\boxed{\text{L3}}
\Longleftrightarrow
\boxed{\text{RAM}}
\]

\subsection{Уровни кэша}

Обычно используются:

\begin{itemize}
    \item \textbf{L1 cache} --- самый быстрый и маленький, часто разделён на кэш команд и кэш данных.
    \item \textbf{L2 cache} --- больше, но медленнее.
    \item \textbf{L3 cache} --- ещё больше, часто общий для нескольких ядер.
\end{itemize}

Схема многоядерного процессора:

\[
\begin{array}{ccc}
\boxed{
\begin{array}{c}
\text{Core 1}\\
\text{L1}\\
\text{L2}
\end{array}}
&
\boxed{
\begin{array}{c}
\text{Core 2}\\
\text{L1}\\
\text{L2}
\end{array}}
&
\boxed{
\begin{array}{c}
\text{Core 3}\\
\text{L1}\\
\text{L2}
\end{array}}
\\[1em]
\multicolumn{3}{c}{
\boxed{\text{Shared L3 Cache}}
}
\\[1em]
\multicolumn{3}{c}{
\boxed{\text{RAM}}
}
\end{array}
\]

\subsection{Кэш-линия}

Данные загружаются в кэш не по одному байту, а блоками фиксированного размера --- \textbf{кэш-линиями}.

Типичный размер кэш-линии:

\[
64 \text{ байта}.
\]

Если программа обращается к адресу $A$, в кэш загружается целый блок:

\[
[A, A+63].
\]

Это эффективно благодаря пространственной локальности.

Пример:

\[
\texttt{int a[1000];}
\]

Если \texttt{int} занимает 4 байта, то в кэш-линию 64 байта помещается:

\[
64 / 4 = 16
\]

элементов массива.

Если загружен \texttt{a[0]}, то, скорее всего, рядом окажутся и:

\[
a[1], a[2], \ldots, a[15].
\]

\subsection{Cache hit и cache miss}

\textbf{Cache hit} --- обращение к данным, которые уже есть в кэше.

\textbf{Cache miss} --- данные отсутствуют в кэше, требуется обращение к более медленному уровню памяти.

Среднее время доступа:

\[
AMAT = T_{hit} + MissRate \cdot MissPenalty.
\]

Где:
\begin{itemize}
    \item $AMAT$ --- Average Memory Access Time;
    \item $T_{hit}$ --- время доступа при попадании;
    \item $MissRate$ --- доля промахов;
    \item $MissPenalty$ --- штраф за промах.
\end{itemize}

Пример:

\[
T_{hit} = 1 \text{ нс},
\quad
MissRate = 5\% = 0.05,
\quad
MissPenalty = 50 \text{ нс}.
\]

Тогда:

\[
AMAT = 1 + 0.05 \cdot 50 = 3.5 \text{ нс}.
\]

\subsection{Организация кэша}

Кэш состоит из строк:

\[
\begin{array}{|c|c|c|}
\hline
\text{Valid bit} & \text{Tag} & \text{Data block} \\
\hline
\end{array}
\]

Адрес памяти разбивается на части:

\[
\boxed{\text{Tag}}
\boxed{\text{Index}}
\boxed{\text{Offset}}
\]

\begin{itemize}
    \item \textbf{Offset} --- смещение внутри кэш-линии;
    \item \textbf{Index} --- номер набора или строки кэша;
    \item \textbf{Tag} --- старшие биты адреса для проверки совпадения.
\end{itemize}

Схема:

\[
\boxed{\text{Address}}
=
\boxed{\text{Tag}}
\boxed{\text{Index}}
\boxed{\text{Offset}}
\]

\[
\boxed{\text{Index}}
\rightarrow
\boxed{\text{Cache Set}}
\]

\[
\boxed{\text{Tag}}
\stackrel{?}{=}
\boxed{\text{Stored Tag}}
\]

Если совпадает и бит valid равен 1, то это cache hit.

\subsection{Прямое отображение}

\textbf{Кэш с прямым отображением} --- каждая строка памяти может попасть только в одно определённое место кэша.

Если в кэше $N$ строк, то номер строки:

\[
index = block\_number \bmod N.
\]

Преимущества:
\begin{itemize}
    \item простая реализация;
    \item быстрый доступ;
    \item дешёвая аппаратура.
\end{itemize}

Недостаток:
\begin{itemize}
    \item возможны частые конфликтные промахи.
\end{itemize}

Пример:

Пусть $N = 4$. Тогда блоки:

\[
0,4,8,12,\ldots
\]

попадают в одну и ту же строку кэша.

\[
\begin{array}{|c|c|}
\hline
\text{Строка кэша} & \text{Блоки памяти} \\
\hline
0 & 0,4,8,12,\ldots \\
\hline
1 & 1,5,9,13,\ldots \\
\hline
2 & 2,6,10,14,\ldots \\
\hline
3 & 3,7,11,15,\ldots \\
\hline
\end{array}
\]

\subsection{Полностью ассоциативный кэш}

В полностью ассоциативном кэше любой блок памяти может быть помещён в любую строку кэша.

Преимущества:
\begin{itemize}
    \item минимальны конфликтные промахи;
    \item гибкое размещение.
\end{itemize}

Недостатки:
\begin{itemize}
    \item сложная аппаратура поиска;
    \item нужно сравнивать тег со всеми строками;
    \item дорого при большом размере.
\end{itemize}

\subsection{Множественно-ассоциативный кэш}

Компромиссный вариант --- \textbf{множественно-ассоциативный кэш}.

Кэш делится на наборы, каждый набор содержит несколько строк.

Если кэш $k$-ассоциативный, то каждый блок может попасть в один набор, но в любую из $k$ строк этого набора.

Например, 4-way set associative cache:

\[
\begin{array}{|c|c|c|c|c|}
\hline
\text{Set 0} & Way 0 & Way 1 & Way 2 & Way 3 \\
\hline
\text{Set 1} & Way 0 & Way 1 & Way 2 & Way 3 \\
\hline
\text{Set 2} & Way 0 & Way 1 & Way 2 & Way 3 \\
\hline
\end{array}
\]

Адрес:

\[
\boxed{\text{Tag}}
\boxed{\text{Set Index}}
\boxed{\text{Offset}}
\]

\subsection{Политики замещения}

Если нужно загрузить новый блок в заполненный набор, надо выбрать, какой блок вытеснить.

Основные политики:
\begin{itemize}
    \item \textbf{LRU} --- Least Recently Used, вытеснить давно не использовавшийся блок;
    \item \textbf{FIFO} --- вытеснить блок, загруженный раньше всех;
    \item \textbf{Random} --- случайный выбор;
    \item \textbf{Pseudo-LRU} --- приближённый LRU, часто используется в железе.
\end{itemize}

\subsection{Политики записи}

При записи данных в кэш возможны разные стратегии.

\subsubsection{Write-through}

\textbf{Write-through} --- данные записываются одновременно в кэш и в следующий уровень памяти.

\[
\boxed{\text{CPU}}
\rightarrow
\boxed{\text{Cache}}
\rightarrow
\boxed{\text{RAM}}
\]

Преимущества:
\begin{itemize}
    \item память всегда содержит актуальные данные;
    \item проще обеспечивать согласованность.
\end{itemize}

Недостатки:
\begin{itemize}
    \item больше трафик в память;
    \item ниже производительность записи.
\end{itemize}

\subsubsection{Write-back}

\textbf{Write-back} --- данные сначала изменяются только в кэше, а в память записываются позже, при вытеснении строки.

Для этого используется \textbf{dirty bit} --- бит модификации.

\[
\boxed{\text{CPU}}
\rightarrow
\boxed{\text{Cache}}
\quad
\text{а RAM обновляется позже}
\]

Преимущества:
\begin{itemize}
    \item меньше обращений к памяти;
    \item выше производительность.
\end{itemize}

Недостатки:
\begin{itemize}
    \item сложнее логика;
    \item нужна поддержка согласованности;
    \item при сбое данные могут быть только в кэше.
\end{itemize}

\subsection{Когерентность кэшей}

В многоядерных системах каждое ядро может иметь собственный кэш. Возникает проблема: разные кэши могут хранить копии одной и той же области памяти.

Пример:

\[
\begin{array}{ccc}
\boxed{\text{Core 1 Cache}} & & \boxed{\text{Core 2 Cache}} \\
x = 5 & & x = 5
\end{array}
\]

Если Core 1 меняет:

\[
x := 10,
\]

то Core 2 не должен продолжать бесконечно видеть старое значение:

\[
x = 5.
\]

Для решения используются протоколы когерентности, например MESI.

Состояния MESI:
\begin{itemize}
    \item \textbf{Modified} --- строка изменена и отличается от памяти;
    \item \textbf{Exclusive} --- строка есть только в данном кэше и совпадает с памятью;
    \item \textbf{Shared} --- строка может быть в нескольких кэшах;
    \item \textbf{Invalid} --- строка недействительна.
\end{itemize}

\section{Современные многоядерные и многопроцессорные системы}

\subsection{Многоядерный процессор}

\textbf{Многоядерный процессор} содержит несколько вычислительных ядер на одном кристалле.

\[
\boxed{
\begin{array}{ccc}
\boxed{\text{Core 1}} & \boxed{\text{Core 2}} & \boxed{\text{Core 3}} \\
\multicolumn{3}{c}{\boxed{\text{Shared L3 Cache}}} \\
\multicolumn{3}{c}{\boxed{\text{Memory Controller}}}
\end{array}}
\]

Преимущества:
\begin{itemize}
    \item параллельное выполнение потоков;
    \item повышение производительности;
    \item эффективное использование энергии;
    \item поддержка многозадачности.
\end{itemize}

\subsection{SMP}

\textbf{SMP} --- Symmetric Multiprocessing, симметричная многопроцессорная система.

Особенности:
\begin{itemize}
    \item несколько процессоров имеют равноправный доступ к общей памяти;
    \item одна операционная система управляет всеми процессорами;
    \item любой поток может выполняться на любом процессоре.
\end{itemize}

Схема:

\[
\begin{array}{ccc}
\boxed{\text{CPU 1}} & \boxed{\text{CPU 2}} & \boxed{\text{CPU 3}} \\
\multicolumn{3}{c}{\updownarrow} \\
\multicolumn{3}{c}{\boxed{\text{Shared Memory}}}
\end{array}
\]

\subsection{NUMA}

\textbf{NUMA} --- Non-Uniform Memory Access. В такой системе память физически распределена между процессорными узлами, но логически образует общее адресное пространство.

\[
\begin{array}{ccc}
\boxed{
\begin{array}{c}
\text{CPU 1}\\
\text{Local RAM 1}
\end{array}}
&
\Longleftrightarrow
&
\boxed{
\begin{array}{c}
\text{CPU 2}\\
\text{Local RAM 2}
\end{array}}
\end{array}
\]

Доступ к локальной памяти быстрее, чем к удалённой.

\[
T_{local} < T_{remote}.
\]

Операционная система старается размещать данные ближе к процессору, который их использует.

\section{Пример полного пути выполнения команды}

Рассмотрим команду:

\[
\texttt{ADD R1, [X]}
\]

Смысл:

\[
R1 := R1 + Memory[X].
\]

Путь выполнения:

\begin{enumerate}
    \item Счётчик команд PC содержит виртуальный адрес команды.
    \item MMU преобразует виртуальный адрес команды в физический.
    \item Проверяется TLB.
    \item Команда извлекается из L1 instruction cache.
    \item Если её нет в L1, поиск идёт в L2, L3, RAM.
    \item Команда декодируется.
    \item Для адреса \texttt{X} выполняется трансляция виртуального адреса данных.
    \item Проверяется TLB данных.
    \item Данные ищутся в L1 data cache.
    \item Если данные отсутствуют, они загружаются из L2/L3/RAM.
    \item АЛУ складывает значение из памяти с регистром \texttt{R1}.
    \item Результат записывается в \texttt{R1}.
\end{enumerate}

Схема:

\[
\boxed{\text{PC}}
\rightarrow
\boxed{\text{TLB}}
\rightarrow
\boxed{\text{I-cache}}
\rightarrow
\boxed{\text{Decode}}
\rightarrow
\boxed{\text{D-TLB}}
\rightarrow
\boxed{\text{D-cache}}
\rightarrow
\boxed{\text{ALU}}
\rightarrow
\boxed{\text{Register}}
\]

\section{Итог}

Современный компьютер представляет собой сложную иерархическую систему, в которой:
\begin{itemize}
    \item процессор выполняет команды программы;
    \item память организована в виде иерархии от регистров до внешних накопителей;
    \item шины и межсоединения обеспечивают обмен данными между устройствами;
    \item виртуальная память создаёт для процессов удобное и защищённое адресное пространство;
    \item страничная организация позволяет эффективно отображать виртуальные адреса на физические;
    \item сегментная организация отражает логическую структуру программы;
    \item кэш-память уменьшает среднее время доступа к данным;
    \item современные процессоры используют конвейеры, кэши, TLB, предсказание переходов, суперскалярность и многоядерность.
\end{itemize}

Краткая обобщающая схема:

\[
\begin{array}{c}
\boxed{
\begin{array}{c}
\text{Программа} \\
\downarrow \\
\text{Виртуальные адреса} \\
\downarrow \\
\text{MMU + TLB} \\
\downarrow \\
\text{Физические адреса} \\
\downarrow \\
\text{Кэш L1/L2/L3} \\
\downarrow \\
\text{RAM} \\
\downarrow \\
\text{SSD/HDD}
\end{array}}
\end{array}
\]

Главная идея современной архитектуры --- скрыть медленность памяти и устройств ввода-вывода за счёт иерархий, кэширования, параллелизма и аппаратной поддержки операционной системы.

\end{document}
